%!TEX root = ../main.tex

% ============================================================
%  CAPÍTULO 1: INTRODUCCIÓN
% ============================================================

\chapter{Introducción}\label{ch:introduccion}

% ──────────────────────────────────────────────────────────────
\section{Motivación del problema}
% ──────────────────────────────────────────────────────────────


La industria global de investigación de mercados y opinión pública, que incluye el pronóstico electoral, alcanzó 153.000 millones de dólares en 2024, un crecimiento del 8\% frente al año anterior, según \textit{ESOMAR} \parencite{esomar2025}, la asociación mundial del sector. En Colombia, la misma industria facturó 743.058 millones de pesos en 2024 (equivalentes a aproximadamente 208 millones de dólares a la tasa de cambio de mayo de 2026), con un crecimiento del 1\,\% frente al año anterior, según la \textit{Asociación Colombiana de Empresas de Investigación de Mercados y Opinión Pública} \parencite{acei2024}. 

En paralelo, el auge de los mercados de predicción financiera ilustra que existe un apetito de mercado creciente por pronósticos ágiles y verificables sobre eventos futuros. Plataformas como \textit{Polymarket} y \textit{Kalshi}, que no son competidoras directas de las encuestadoras sino mercados financieros de contratos sobre resultados, proyectan alcanzar un billón de dólares en volumen transado anual para 2030, según estimaciones de Bernstein Research \parencite{bernstein2026}. Aunque se trata de industrias de naturaleza distinta, este fenómeno revela una demanda no atendida por información predictiva en tiempo real que las encuestas tradicionales, por diseño, no están estructuradas para satisfacer.

A pesar de su ubicuidad, las encuestas electorales en Colombia enfrentan profundos retos estructurales que se han agudizado recientemente. La dificultad logística de abarcar territorios extensos, la alta tasa de no respuesta y los elevados costos operativos limitan la frecuencia y el tamaño muestral de los estudios, especialmente en elecciones locales.

A estos problemas estructurales se suma un nuevo marco regulatorio que ha tensionado aún más el sector: la \textbf{Ley 2494 de 2025}, que estableció requisitos metodológicos estrictos para la elaboración y publicación de encuestas electorales, incluyendo una veda de publicación que interrumpió el flujo de información entre julio y noviembre de 2025 \parencite{ley2494_2025}. 

La aplicación práctica de esta norma generó tal nivel de fricción entre el \textit{Consejo Nacional Electoral} (CNE) y las firmas encuestadoras, que la consultora española GAD3 suspendió definitivamente sus mediciones en el país en mayo de 2026, argumentando que la interpretación de la Comisión Técnica exige demostrar condiciones de muestreo que ``ninguna metodología demoscópica puede satisfacer en la práctica'' \parencite{gad3_2026}. 

La norma, al invalidar encuestas telefónicas, paneles digitales y métodos mixtos en favor de encuestas presenciales domiciliarias, encareció operativamente el sector y concentró el mercado en pocas firmas con suficiente capacidad económica para operar en campo en todo el territorio nacional \parencite{occidente2026}. En este panorama, resulta cada vez más urgente desarrollar herramientas de pronóstico electoral que no dependan exclusivamente de encuestas y que aprovechen fuentes alternativas de datos masivos.

% ──────────────────────────────────────────────────────────────
\section{Estado del arte: de los modelos paramétricos al SOMEN-DC}
% ──────────────────────────────────────────────────────────────

La capacidad de anticipar resultados electorales ha sido un reto persistente
para la ciencia política desde los años setenta, cuando surgieron los primeros
modelos predictivos formales en Estados Unidos y el Reino Unido. Los trabajos
pioneros de \textcite{sigelman1979} y \textcite{whiteley1979} demostraron que
la popularidad presidencial medida meses antes de una elección era un predictor
estadísticamente significativo del resultado electoral, abriendo la puerta a
una tradición de modelos formales de pronóstico que hasta entonces era
prácticamente inexistente en la ciencia política.

Esta tradición evolucionó rápidamente hacia arquitecturas econométricas más
sofisticadas. \textcite{mughan1987} comparó tres modelos simples de pronóstico
para elecciones generales en el Reino Unido, concluyendo que las variables
macroeconómicas, en particular el desempleo y la aprobación gubernamental,
ofrecían mayor poder predictivo que los indicadores de popularidad aislados.
Por su parte, \textcite{sanders1991} desarrolló un modelo de función de
popularidad para el gobierno británico que incorporaba condiciones económicas
y eventos políticos exógenos, anticipando el resultado de las elecciones
generales con seis meses de antelación. Durante los años noventa y la primera
década del siglo XXI, estos modelos alcanzaron su madurez técnica,
apoyándose en herramientas de series de tiempo como VAR y ARMA. En el
contexto del Reino Unido, \textcite{lewisbeck2004} formularon un modelo de
economía política que explicaba más del 80\,\% de la varianza en el apoyo
gubernamental usando solo tres variables: aprobación del gobierno, inflación
y el número de términos consecutivos en el poder. En paralelo,
\textcite{norpoth2004} propuso una perspectiva dinámica que modelaba el retorno
del voto al equilibrio histórico mediante un proceso autorregresivo AR(1),
demostrando que el partido ganador tiende a perder aproximadamente la mitad
de su ventaja en la siguiente elección. En varios ciclos electorales, estos
modelos superaban en precisión a las propias encuestadoras.

Sin embargo, desde aproximadamente 2015, la capacidad predictiva de estos
enfoques se deterioró de forma pronunciada. En el Reino Unido, ningún modelo
establecido logró anticipar la mayoría conservadora de ese año
\parencite{stegmaier2023}; un año después, los mismos enfoques fallaron en
predecir la victoria de Donald Trump, el Brexit y los resultados de otras consultas populares de ese ciclo \parencite{stegmaier2023, brito2024stop}. Esta crisis
generalizada no fue accidental: \textcite{kennedy2018} documentaron, en una
evaluación exhaustiva encargada por la Asociación Americana de Investigación
de Opinión Pública (AAPOR), que los sondeos de 2016 sobreestimaron
sistemáticamente el apoyo a Hillary Clinton debido a sesgos de no-respuesta
diferencial, especialmente entre votantes de Trump sin estudios universitarios.
A esto se sumaron los cambios estructurales en la comunicación política
propiciados por las redes sociales \parencite{tucker2018} y la creciente
fragmentación y polarización de los electorados \parencite{barbera2020}, que
volvieron menos predecibles los comportamientos que los modelos paramétricos
habían aprendido a capturar.

La crisis impulsó a los investigadores a explorar fuentes de datos
alternativas. El auge del \textit{Big Data} y las plataformas digitales abrió
una nueva línea de investigación: usar las huellas digitales de los ciudadanos
como covariables en tiempo real para el pronóstico electoral. El estudio
seminal de \textcite{tumasjan2010} sobre las elecciones federales alemanas de
2009 fue el primer trabajo de referencia en este campo. Analizando más de
104.000 \textit{tweets} publicados en las semanas previas a la elección,
los autores encontraron que el simple volumen de menciones de cada partido
replicaba el resultado electoral con un error absoluto medio de apenas
1{,}65\,\%, comparable al de las encuestadoras tradicionales. Más allá de
la predicción, Tumasjan et al.\ demostraron que los perfiles de sentimiento
de los mensajes correspondían a las posiciones políticas reales de los
partidos, y que las menciones conjuntas de dos partidos reflejaban las
alianzas políticas existentes. Esta hipótesis, que la atención digital
equivale a capital político, abrió una década de investigación intensa.

No obstante, la promesa inicial no se sostuvo en la replicación sistemática.
\textcite{skoric2020}, en un meta-análisis de 74 estudios publicados entre
2007 y 2018, encontraron que las predicciones basadas en redes sociales
quedaban en promedio por debajo de los \textit{benchmarks} de la encuesta
tradicional, con alta varianza entre estudios y sin convergencia metodológica.
Sus resultados mostraron que las predicciones basadas únicamente en análisis
de sentimiento léxico no superaban a las basadas en características
estructurales de red, y que la combinación de ambas, usando aprendizaje
automático, producía los mejores resultados (MAE = 2{,}14\,\%;
$R^2 = 0{,}818$). \textcite{brito2024stop}, en una revisión crítica de doce
años de investigación, formalizaron el problema: el enfoque de
\textit{Twitter} solo, basado en volumen de menciones o sentimiento, no rinde
mejor que el azar, con tasas de acierto de entre el 50\,\% y el 63\,\%.
Los autores identificaron cinco causas estructurales de este fracaso: la
imposibilidad de encuestar representativamente a través de Twitter, la
dependencia de decisiones arbitrarias de recolección, los cambios constantes
en el ecosistema de plataformas, la falta de datos históricos comparables,
y la repetición de los errores metodológicos de las \textit{straw polls}
anteriores a 1936. El problema de fondo es metodológico: rastrear millones
de publicaciones abiertas es costoso, ruidoso y difícilmente replicable
fuera de entornos universitarios con infraestructura computacional avanzada
\parencite{skoric2020}.

La respuesta más sofisticada a estos problemas llegó desde Brasil.
\textcite{brito2020predicting} propusieron un cambio fundamental de paradigma:
en lugar de rastrear lo que la ciudadanía dice sobre los candidatos,
menciones abiertas, \textit{hashtags}, sentimiento, el modelo analiza
exclusivamente lo que la ciudadanía hace con el contenido que los propios
candidatos publican en sus perfiles oficiales. Combinando datos de Facebook,
Twitter e Instagram de las elecciones presidenciales de Brasil (2018) y
Estados Unidos (2016) con encuestas tradicionales como variable objetivo,
entrenaron redes neuronales tipo MLP que superaron estadísticamente a las
encuestadoras en Brasil, donde los candidatos son muchos y las encuestas
son pocas, y fueron competitivos en Estados Unidos. Este enfoque fue
posteriormente sistematizado por \textcite{brito2022measuring}, quienes
analizaron más de 67.000 publicaciones de candidatos presidenciales en
Argentina, Brasil, Colombia y México (2018--2019), demostrando que las
métricas de \textit{engagement} por publicación, en particular los
\textit{likes} por \textit{post} en Instagram, correlacionaban fuertemente
con el resultado electoral ($r = 0{,}75$--$0{,}98$), mientras que el
número absoluto de publicaciones presentaba correlación nula o negativa.

El conjunto de estas contribuciones se formalizó en el marco
\textbf{SOMEN-DC} (\textit{Social Media Electoral Nowcasting During Campaign}),
desarrollado por \textcite{brito2023somen} en la Universidad Federal de
Pernambuco. El modelo propone un proceso de cinco pasos: comprensión del evento, recolección de datos, preparación, modelado y evaluación, y una
estrategia de ejecución continua (DC) que permite emitir predicciones diarias
durante la campaña. Usando datos de más de 65.000 publicaciones y 195
encuestas de cuatro países latinoamericanos, los ensambles de redes neuronales
(MLP y GRNN) entrenados con métricas de \textit{engagement} y encuestas como
etiquetas alcanzaron errores absolutos medios frecuentemente inferiores a
3 puntos porcentuales, dentro del umbral histórico de las encuestadoras.

Sin embargo, el \textsc{somen-dc} presenta tres limitaciones estructurales que
este trabajo busca superar:

\begin{enumerate}
  \item \textbf{Techo encuestocéntrico.} Al usar las encuestas como variable
        objetivo de entrenamiento, el modelo no puede superar la capacidad
        predictiva de las propias encuestas. Es, en esencia, un replicador
        sofisticado de sondeos: aprende a imitar la señal de las encuestas
        a partir de datos digitales, sin acceso a una fuente de verdad
        independiente.

  \item \textbf{Dependencia de infraestructura costosa.} El modelo requiere
        \textit{scraping} continuo desde el inicio de la campaña para registrar
        la evolución temporal de las interacciones, acceso que los autores
        originales tenían a través de servicios institucionales de la
        Universidad Federal de Pernambuco, hoy inasequibles o inexistentes
        para la mayoría de investigadores independientes.

  \item \textbf{No transferibilidad.} Los pesos aprendidos son válidos
        exclusivamente para la elección en que se entrenaron. Cada nueva
        contienda exige repetir desde cero el ciclo completo de recolección y
        entrenamiento, sin posibilidad de acumular aprendizaje entre elecciones
        o de generalizar patrones entre contextos electorales distintos.
\end{enumerate}

Estas tres limitaciones definen el espacio de contribución de esta tesis.

% ──────────────────────────────────────────────────────────────
\section{Hipótesis de trabajo}
% ──────────────────────────────────────────────────────────────

El punto de partida conceptual de este proyecto es la hipótesis seminal
formulada por \textcite{tumasjan2010}, quienes postulan que el volumen
de actividad digital de un candidato en plataformas de redes sociales
funciona como un estimador de su capital electoral subyacente. Esta
intuición ha encontrado respaldo empírico consistente en contextos
distintos al original: \textcite{brito2022measuring} demostraron que
las métricas de \textit{engagement} por publicación en los perfiles
oficiales de candidatos presidenciales en Argentina, Brasil, Colombia
y México correlacionan fuertemente con el resultado electoral
($r = 0{,}75$--$0{,}98$), mientras que las interacciones absolutas o
el volumen de publicaciones presentan correlaciones nulas o negativas.

No obstante, existe un argumento estructural en contra de esta
hipótesis que merece atención explícita: la penetración de las redes
sociales no es aleatoria. La probabilidad de que un ciudadano tenga
presencia activa en redes sociales depende de factores económicos,
geográficos y culturales, lo que introduce un sesgo sistemático si se
asume que la audiencia digital de un candidato representa fielmente
a su electorado potencial \parencite{skoric2020}. Sin embargo, este
argumento pierde fuerza en la medida en que la cobertura digital se
aproxima a la universalidad. Según \textcite{datareportal2026}, el
70{,}4\,\% de la población colombiana usa redes sociales al cierre de
2025, cifra que creció en 1{,}6 millones de personas solo en 2024,
con Facebook, TikTok, YouTube, Instagram y LinkedIn alcanzando al
89{,}3\,\%, 93{,}0\,\%, 74{,}4\,\%, 52{,}3\,\% y 44{,}4\,\% de los
usuarios de redes sociales colombianos, respectivamente (porcentajes
calculados sobre la base de quienes ya usan redes sociales, no sobre
la población adulta total). Significativamente, X/Twitter, la plataforma que concentra la mayor parte de la literatura de
predicción electoral con redes sociales, solo alcanza al 12{,}0\,\%
de los usuarios de redes sociales en Colombia, lo que compromete estructuralmente
cualquier modelo que dependa exclusivamente de esa fuente.

\subsection*{Formulación del problema}

Considérese una elección de ganador único para un cargo ejecutivo: alcalde, gobernador o presidente, con un conjunto finito de
$C$ candidatos competitivos $\mathcal{C} = \{c_1, c_2, \ldots, c_C\}$.
Para cada candidato $c_i$, se observa un vector de características
digitales:

\begin{equation}
    \mathbf{x}_i \in \mathbb{R}^d
    \label{eq:vector-caracteristicas}
\end{equation}

\noindent construido a partir de las interacciones acumuladas en sus
publicaciones de redes sociales durante una ventana de $n$ días previos
a la jornada electoral. Este vector puede incluir interacciones totales,
\textit{engagement} por publicación, dominancia relativa respecto al
total de interacciones de la elección, entre otras métricas.

La hipótesis central de este trabajo postula que existe una función:

\begin{equation}
    f: \mathbb{R}^d \rightarrow [0, 1]
    \label{eq:funcion-modelo}
\end{equation}

\noindent tal que $f(\mathbf{x}_i)$ aproxima la cuota de voto esperada
del candidato $c_i$. La restricción de coherencia, que las predicciones
sumen a la unidad sobre todos los candidatos de la misma elección, no
es una propiedad de $f$ en sí misma, sino una normalización que se impone
a posteriori sobre el conjunto de salidas escalares:

\begin{equation}
    \sum_{i=1}^{C} f(\mathbf{x}_i) = 1
    \label{eq:restriccion-suma}
\end{equation}

La contribución metodológica de este trabajo respecto a la literatura
previa reside en el mecanismo de aprendizaje de $f$. En los modelos
existentes, incluyendo el \textsc{somen-dc}
\parencite{brito2023somen}, esta función se aprende de forma
independiente para cada elección, usando las encuestas de intención de
voto como variable objetivo. En el modelo propuesto, $f$ se aprende
sobre un conjunto de $M$ elecciones históricas
$\mathcal{E} = \{e_1, e_2, \ldots, e_M\}$, usando directamente el
resultado electoral como fuente de verdad. Esto tiene dos implicaciones
fundamentales: primero, elimina la dependencia de encuestas en cualquier
etapa del proceso; segundo, permite que los parámetros aprendidos
capturen patrones que son comunes entre elecciones, la relación
estructural entre tracción digital y respaldo electoral, y no los
detalles particulares de una campaña en específico. Una vez entrenada,
$f$ puede aplicarse a una elección nueva $e_{M+1} \notin \mathcal{E}$
sin necesidad de datos de entrenamiento adicionales, lo que constituye
la propiedad de \textbf{elección-agnosticismo} que da nombre al modelo.

Bajo este marco, el modelo opera recibiendo como \textit{input} el
vector $\mathbf{x}_i$ de interacciones digitales de cada candidato,
calculado sobre una ventana de $n$ días previos a la jornada electoral.
El \textit{output} se obtiene aplicando $f$ individualmente a cada candidato:
$\hat{y}_i = f(\mathbf{x}_i) \in [0, 1]$, y ensamblando el vector resultante
$\hat{\mathbf{y}} = (\hat{y}_1, \ldots, \hat{y}_C) \in [0,1]^C$, donde cada
componente representa la cuota de voto estimada para el candidato $c_i$,
respondiendo a la pregunta: \textit{si las elecciones ocurrieran hoy, \textquestiondown{}qué porcentaje de los votantes apoyaría a cada candidato?}

% ──────────────────────────────────────────────────────────────
\section{Estructura del documento}
% ──────────────────────────────────────────────────────────────

Para abordar esta hipótesis, la investigación se estructuró a través de tres etapas empíricas que guían la organización de este documento. El \textbf{Capítulo 2} define formalmente la pregunta de investigación y los interrogantes derivados. El \textbf{Capítulo 3} presenta el marco conceptual, explorando la literatura sobre \textit{nowcasting} electoral y la viabilidad de los modelos elección-agnósticos. Seguidamente, el \textbf{Capítulo 4} detalla los objetivos del estudio.

El \textbf{Capítulo 5} desarrolla la metodología, dividida en las tres fases del proyecto: 
\begin{itemize}
    \item \textbf{Bolivia (Prueba de concepto):} donde se replica el modelo SOMEN-DC con herramientas comerciales, documentando sus limitaciones metodológicas.
    \item \textbf{Costa Rica (Modelado de crecimiento):} enfocada en resolver el problema del registro histórico mediante la calibración matemática de la curva de crecimiento de las interacciones.
    \item \textbf{Colombia (Construcción y modelado final):} donde se compila el \textit{dataset} elección-agnóstico de elecciones locales (2023) y se entrenan los modelos de regresión fraccional (evaluando la precisión de clasificación del \textit{Top-1} derivado).
\end{itemize}

Finalmente, el \textbf{Capítulo 6} expone los resultados obtenidos en cada fase, mientras que los \textbf{Capítulos 7 y 8} discuten el impacto esperado de la herramienta, sus limitaciones inherentes y las conclusiones definitivas del estudio.
